\chapter{The Tube: BBC BASIC}\label{cha:tube}

The BBC Micro's most elegant idea was the Tube: a fast port through which a
\emph{second processor} --- another CPU with its own memory --- could take
over the computation while the Beeb kept the keyboard, the screen and the
discs. The BMC-K4510 has a Tube of its own at \texttt{\$D800}, and the first
thing fitted to it is Richard Russell's BBC BASIC, running on the host
machine with a flat 256\,MB of its own.

\begin{type}
BBC
\end{type}

You get the \verb|>| prompt of a BASIC with real power behind it:
\verb|HIMEM=PAGE+250*1024*1024| works, a 200\,MB \verb|DIM| works, and it
is \emph{fast} --- the co-processor is not cycle-matched to 1985. Type
\verb|*QUIT| to hand the console back to the shell.

\section{Graphics and sound}
The console cannot show what it cannot see, so the Tube speaks up: MODE,
GCOL, PLOT, MOVE, DRAW, CIRCLE and the palette arrive at the machine as
escape sequences, and the Tube's gatekeeper (the \emph{Tube ULA}) executes
them on the VICKe blitter. \verb|SOUND| and \verb|ENVELOPE|-less music go
the same way, onto the machine's four-channel sound sequencer and out
through the SIDs --- the classic Beeb \texttt{SOUND chan,amp,pitch,dur},
queued per channel like the original.

\begin{type}
LOAD "BBCBASIC/KALEID.BBC"
RUN
\end{type}

\screen{shots/bbc.png}{\texttt{KALEID.BBC}: MODE 2 kaleidoscope, drawn by
BBC BASIC on the Tube, painted by VICKe.}

The \texttt{/BBCBASIC} directory carries a period programme selection:
\texttt{KALEID}, \texttt{CIRCLES}, \texttt{ROSES}, \texttt{MOUNTAIN},
\texttt{BOUNCE}, \texttt{CLOCK} (a live analogue clock face) and
\texttt{TUNE} --- Fr\`ere Jacques as a three-voice round, the sound
sequencer's party piece.

\section{Star commands}
A line starting with \verb|*| goes to the machine: \verb|*DIR|, \verb|*CD|,
any shell command --- and BBC BASIC's own file words (\verb|LOAD|,
\verb|SAVE|) read and write the machine's filesystem directly, because the
co-processor lives inside \texttt{fs/} too.

\begin{aside}
\textbf{Edges, honestly:} \texttt{POINT(} and \texttt{TINT} are not
implemented; \texttt{GCOL} modes 1--4 draw plain; \texttt{VDU 5} text
prints as text. The Tube runs on the desktop host only for now --- the Pi
build does not fit a co-processor yet.
\end{aside}
